% !TEX root = ../glossario.tex
% Lettera D del Glossario (vedi capitoli/a.tex per le istruzioni).
\phantomsection
\letteraglossario{D}

\termine{Data warehouse}{Archivio di dati ottimizzato per l'analisi, in cui i dati dei sistemi
  gestionali vengono puliti, integrati e organizzati per essere
  interrogati in modo efficiente, per esempio nella \glo{Business
  Intelligence}.}

\termine{Design}{Fase del ciclo di vita del software in cui i requisiti vengono tradotti in una specifica architettura e struttura del sistema, definendo le componenti, le interfacce e i dati necessari.}

\termine{Design pattern}{Soluzione progettuale generale, riutilizzabile e consolidata per un problema ricorrente all'interno di un determinato contesto nello sviluppo del software.}

\termine{Documentation as Code}{Pratica che consiste nel gestire la stesura, il versionamento, la revisione e la pubblicazione della documentazione utilizzando gli stessi strumenti e processi adottati per il codice sorgente (es. Git, repository, linguaggi di markup, integrazione continua).}

\termine{Dominio applicativo}{Area tematica e operativa in cui il software deve operare.}